\chapter{Literature Review}

\section{Overview of Fake News Detection}

\begin{itemize}
    \item Introduce fake news and inauthentic content problem
    \item Explain importance of detection in social media
    \item Evolution of approaches:
    \begin{itemize}
        \item Text-only models
        \item Context-aware models
        \item Multimodal models
    \end{itemize}
    \item Content alone is not sufficient, context and behaviour matter
    \item Deep learning dominates modern approaches
    \item Different approaches address different aspects of the problem
\end{itemize}

\section{Content-Based Approaches}

\begin{itemize}
    \item Models analyse only textual features
    \item Techniques:
    \begin{itemize}
        \item NLP, BERT, RNN, CNN
    \end{itemize}
    \item Examples:
    \begin{itemize}
        \item Text branch in DANES and GETAE
    \end{itemize}
    \item Strengths:
    \begin{itemize}
        \item Simple and efficient
        \item Good baseline performance
    \end{itemize}
    \item Weaknesses:
    \begin{itemize}
        \item Cannot detect visual misinformation
        \item Ignores user behaviour and propagation
        \item Struggles with subtle misinformation
    \end{itemize}
    \item Content-only approaches are not sufficient for real-world scenarios
\end{itemize}

\section{Context-Aware Approaches}

\begin{itemize}
    \item Use social media context instead of only content
    \item Types of context:
    \begin{itemize}
        \item User behaviour
        \item Interactions
        \item Propagation patterns
    \end{itemize}
    \item Models:
    \begin{itemize}
        \item DANES (text + social context)
        \item GETAE (text + propagation graph)
        \item GCAN (context only)
    \end{itemize}
    \item Techniques:
    \begin{itemize}
        \item Graph neural networks
        \item Propagation graphs
    \end{itemize}
    \item Strengths:
    \begin{itemize}
        \item Captures real-world behaviour
        \item Improves detection performance
    \end{itemize}
    \item Weaknesses:
    \begin{itemize}
        \item Relies on text-only content
        \item Ignores multimodal data
        \item Depends on specific datasets
    \end{itemize}
    \item Context improves detection but content is still limited
\end{itemize}

\section{Multimodal Approaches}

\begin{itemize}
    \item Combine multiple data types such as text and images
    \item Social media content is inherently multimodal
    \item Models:
    \begin{itemize}
        \item EANN, MCNN, HMCAN, MPFN, LDPG
    \end{itemize}
    \item Techniques:
    \begin{itemize}
        \item CNN, transformers, fusion methods
    \end{itemize}
    \item Strengths:
    \begin{itemize}
        \item Better feature representation
        \item Improved accuracy
        \item Detects misleading visuals
    \end{itemize}
    \item Weaknesses:
    \begin{itemize}
        \item Ignores social context
        \item Static models without propagation
        \item Requires large datasets
    \end{itemize}
    \item Multimodal improves content but lacks context integration
\end{itemize}

\section{Research Gap and Positioning}

\begin{itemize}
    \item Content-based models are limited
    \item Context-aware models rely on text-only content
    \item Multimodal models ignore social context
    \item There is no effective integration of multimodal and context-aware approaches
    \item Existing models such as DANES and GETAE use text-only content branches
    \item Proposed research replaces text branch with multimodal branch
    \item Combines multimodal content with context-aware architecture
    \item Evaluates performance improvement and contribution of each branch
\end{itemize}


